\section{Why the Composition Fails}
\label{sec:why}

The scope of what follows is fixed before the argument rather than after it.
This is an analysis of a metric's algebra, illustrated by measurements on two
corpora. Properties (a) and (b) are consequences of the definitions quoted in
Section~\ref{sec:metric} and hold wherever those definitions are used. Property
(c) is a statement about what the algebra implies, not a survey finding.

The defect is not a badly chosen weight. It is an error of type. Write the
score as $B + U$, where $B$ is the weighted sum of AP and AUC and is bounded
above by the sum of their weights, and $U$ is the weighted sum of the timing
terms, bounded above by the sum of their weights times the mean accident onset
--- a quantity that depends on how the corpus was windowed.

\emph{(a) The unbounded term grows with clip length while the bounded one does
not.} The ratio of their maxima scales linearly in the mean onset, so two
groups adopting the same weights on corpora windowed differently are not using
the same metric, though they report the same name. On the competition the
maximum of $U$ is $2.00186$, measured; the maximum of $B$ is at least $0.49979$,
giving a ratio of at most $4.01$, and against $B$ at chance the measured ratio
is $5.70$. Section~\ref{sec:nexar} shows the effect does not even need two
corpora: within one, collisions ranging from $3.03$ to $56.80$\,s give an
18.7-fold spread in the earliness a single constant can collect.

\emph{(b) The unbounded term has a constant as its global maximiser.} Any
\emph{unconditioned} threshold-crossing earliness measure is maximised by
crossing immediately and never returning. The maximiser reads no input, so the
term it maximises can carry no information about the input. The qualifier is
the seed of the remedy: an earliness measure that is conditioned --- evaluated
at a fixed operating point of the discrimination metric, or penalised by the
false-alarm rate --- is not maximised by a constant.

\emph{(c) Consequently the separating power of the score lies in $B$.} To the
extent that a submission crosses early and holds, its $U$ is at or near maximum
and it is separated from other such submissions only by $B$. We cannot verify
this for any particular entry, having no team's curves but our own, so we state
it as what the algebra implies and not as an observation about anyone's system.

Property (c) is why the failure is hard to notice from inside. A leaderboard of
this shape still orders methods usefully, because $B$ still varies. What is
invisible without a control is that the number attached to each of them is
dominated by a term available for free, so the reported score --- the thing
that goes into a table and gets compared across publications --- is mostly a
constant.

Two mechanisms compound this, and both are visible in our own pipeline rather
than anyone else's. \emph{Metric-enforcing post-processing}: an operation that
clamps the risk score above the threshold, followed by a report of
stable-anticipation compliance, has measured only that the clamp executed.
Stated generally, a post-hoc operation that enforces a metric's definition
renders that metric vacuous. \emph{Non-causal timing gain}: an operation that
resamples the curve so that frame $t$ reports a value computed from a later
frame buys time-to-accident from the future, and whatever earliness it buys
cannot be obtained by any system running in real time. A third and milder
mechanism is selection --- reporting the mean anticipation of the best fifty
clips of 1{,}417 is a statement about the selection, not the method.

\subsection{What this costs safety research}

Read as a system rather than as a curve, the maximiser of (b) is a device that
alerts continuously, on every drive, whether or not anything is in front of the
vehicle. Alarm timing is among the things that determine whether a driver
trusts a forward collision warning and acts on it: warnings arriving when the
driver does not expect them reduce that trust
substantially~\cite{abe2006alarm}, and a warning that is always present is the
limiting case. A composite whose earliness half is maximised by permanent alarm
is therefore not merely uninformative about warning capability; on this
dimension it is ordered the wrong way round. Section~\ref{sec:nexar} puts a
number on the cost the metric cannot see: \SI{60}{\minute} of alarm in every
hour of ordinary driving.

A published anticipation figure therefore does not, on its own, establish that
a system anticipates. A reader given only a composite score cannot tell how
much was earned by the method and how much was available to any submission that
crosses early --- and neither can the authors, unless a control was run.

\section{A Corrected Protocol}
\label{sec:protocol}

Each item is traced to a mechanism above. The first is the one we would keep if
we could keep only one. Items 1--6 were proposed before the validation of
Section~\ref{sec:nexar}; items 7 and 8 exist \emph{because} of it, and are the
substantive change this paper makes to its own recommendation.

\begin{enumerate}
\item \textbf{Report a video-blind control with every anticipation result.}
  Compute the full metric on a constant above threshold, a linear ramp and a
  mid-clip step, and put those rows in the results table. If the control lands
  within noise of the method, the metric is not measuring the method. This
  costs one evaluation run, and should be as reflexive as a majority-class
  baseline in classification.
\item \textbf{Bound the earliness term, and condition it.} Bounding ---
  normalising per clip so achieved over achievable lies in $[0,1]$ --- stops
  the term dominating the sum and stops it depending on the windowing. It does
  \emph{not} remove the constant as maximiser, which still scores $1$ and still
  gets it free; Section~\ref{sec:nexar} confirms this directly. Conditioning on
  a fixed operating point is what removes the exploit. If only one is adopted,
  adopt conditioning.
\item \textbf{State the reference point of every timing metric}: the annotated
  collision, or the end of the clip. Reporting the interval to the end of the
  window under the name time-to-accident inflates every result by a
  corpus-dependent constant.
\item \textbf{Report discrimination alongside any timing metric}, never timing
  alone, and state the ratio of positive to negative clips. Without negatives
  there is no false-positive rate, and any rising curve achieves perfect
  recall.
\item \textbf{Report metrics with post-processing disabled as well as
  enabled.} If a number moves when a clamp is removed, the clamp was part of
  the measurement. If a number moves when a resampling stage is removed, check
  whether that stage reads a frame that has not yet arrived.
\item \textbf{Publish the score of a released control alongside the
  leaderboard.} Organisers can do in an afternoon what took us twenty-five
  submissions, and an entrant would then see at a glance how much of their
  score they earned. For competitions that withhold labels, publish the weights
  and the scorer too: withholding labels prevents overfitting, but withholding
  the weights does not prevent the metric being exploited --- this paper
  exploited it comprehensively without them --- while it does prevent the
  sanity check that would have surfaced the problem before the competition
  opened.
\item \textbf{Average earliness over all positives, not over detected ones.}
  Averaging only where the alarm fired lets a submission buy a perfect score by
  detecting almost nothing; in Section~\ref{sec:nexar} a control scored a
  perfect $1.0000$ on 14 positives of 150. Earliness never delivered is not
  earliness.
\item \textbf{Measure the alarm at frame level, not clip level.} Clip-level
  conditioning cannot see that a curve is flat \emph{within} a clip, and in
  Section~\ref{sec:nexar} a control reading only the video file's header beat a
  model reading the road under items 1--6 alone. Report \emph{coverage} --- the
  share of the annotated actionable window during which the alarm is raised ---
  against \emph{duty cycle} --- the share of all non-actionable time during
  which it is raised --- and summarise as coverage at a capped duty. Report a
  \textbf{metadata-only control} beside the video-blind one, because corpus
  metadata can itself carry the label.
\end{enumerate}

\section{Limitations}
\label{sec:limits}

The leaderboard measurements rest on one competition and one corpus, and the
scorer is closed, so Section~\ref{sec:anomaly} records a discrepancy we could
not resolve. The validation corpus is a second, independent dataset, but it is
still one dataset, and its label leakage through file metadata
(Section~\ref{sec:nexar}) is a property of that corpus whose prevalence
elsewhere we have not measured. The sighted baseline is deliberately weak; it
establishes that the repaired statistic separates sighted from blind, not where
a strong method would fall on it. The actionable window we score against is the
annotators' judgement, and a different annotation convention would move the
coverage numbers though not the ordering. Finally, our claim about the absence
of video-blind controls in the literature covers the works whose results tables
we inspected directly, and is not a systematic audit.

\section{Conclusion}
\label{sec:conclusion}

A composite that adds a bounded discrimination term to an unbounded earliness
term does not measure warning capability. On a live benchmark a constant risk
score reading no pixels matched the eighth of thirteen teams, and the timing
terms were worth $5.70$ times what discrimination contributed to it. On an
independent annotated corpus the same constant took 99.9\% of the available
earliness at exactly chance discrimination, while the only submission that
discriminated earned $0.045$ of $18.86$ seconds.

Three things are consequently cheap that were not. A video-blind control costs
one evaluation run and settles whether a reported score is measuring a method
or a metric. Probe submissions turn a leaderboard into a means of auditing a
closed scorer, without labels and without the weights. And an earliness term
that is bounded, conditioned on a fixed operating point, and measured at frame
level against the time the alarm spends raised gives the field a quantity that
can be compared across corpora windowed differently.

We would stress what the validation changed. The protocol we proposed was
tested where ground truth exists, behaved exactly as predicted on the control
it was designed against, and was then beaten by a control that never decodes a
frame. A protocol that has not been run against adversarial controls is a
hypothesis. What is at stake in getting this right is not tidiness: a warning
system earns its place in a vehicle by alerting a driver early enough to act
and seldom enough to be believed, and a metric whose earliness half is
maximised by a permanent alarm cannot distinguish a system that does that from
one that does nothing. Until anticipation results are reported beside a control
that reads no video, a published warning horizon should be read as an upper
bound on what the method might contribute, not as evidence of what it does.
